\documentclass[12pt]{article}
\usepackage[utf8]{inputenc}
\usepackage[T1]{fontenc}
\usepackage{expl3, xparse}
\usepackage{longtable}
\usepackage[table]{xcolor}

\title{A Multi-Pass Self-Evolving Linguistic System in \LaTeX}
\author{Nikola Topalov}
\date{}

\ExplSyntaxOn

\tl_new:N \l__evo_current_tl
\tl_new:N \l__evo_next_tl
\tl_new:N \g__evo_table_tl
\int_new:N \l__evo_gen_int

\cs_new_protected:Nn \__evo_stage_one:
{
    \tl_replace_all:Nnn \l__evo_current_tl { system } { construct }
    \tl_replace_all:Nnn \l__evo_current_tl { data } { signal }
}

\cs_new_protected:Nn \__evo_stage_two:
{
    \tl_replace_all:Nnn \l__evo_current_tl { the } { THE }
    \tl_put_right:Nn \l__evo_current_tl { ~[expanded] }
}

\cs_new_protected:Nn \__evo_stage_three:
{
    \tl_reverse:N \l__evo_current_tl
    \tl_put_left:Nn \l__evo_current_tl { [CHAOS]~ }
}

\cs_new_protected:Nn \__evo_apply_stage:
{
    \int_case:nn { \l__evo_gen_int }
    {
        {1} { \__evo_stage_one: }
        {2} { \__evo_stage_two: }
        {3} { \__evo_stage_three: }
    }
}

\cs_new_protected:Nn \__evo_run:n
{
    \tl_set:Nn \l__evo_current_tl { #1 }

    \int_step_inline:nn { 3 }
    {
        \int_set:Nn \l__evo_gen_int { ##1 }

        \__evo_apply_stage:

        \tl_gput_right:Nx \g__evo_table_tl
        {
            \int_use:N \l__evo_gen_int & 
            \exp_not:V \l__evo_current_tl
            \exp_not:n { \\ \hline }
        }
    }
}

\NewDocumentCommand{\EvolutionTable}{m}
{
    \tl_gclear:N \g__evo_table_tl

    \begin{longtable}{|c|p{10cm}|}
    \hline
    \rowcolor[gray]{0.9} Generation & Output \\ \hline

    \__evo_run:n { #1 }

    \g__evo_table_tl

    \end{longtable}
}

\ExplSyntaxOff

\begin{document}

\maketitle

\section*{Abstract}
We introduce a self-referential linguistic system implemented entirely in \LaTeX. The system evolves input text through multiple transformation stages, simulating temporal progression and structural mutation.

\section*{Initial Input}
\texttt{The system processes data efficiently}

\section*{Evolution Timeline}

\EvolutionTable{The system processes data efficiently}

\section*{Interpretation}
\begin{itemize}
\item Generation 1 introduces semantic mutation.
\item Generation 2 amplifies structural features.
\item Generation 3 destabilizes the system into nonlinear output.
\end{itemize}

\section*{Conclusion}
This experiment demonstrates that \LaTeX\ can simulate iterative transformation systems with state persistence across generations. The resulting structure resembles a constrained evolutionary process operating entirely within a typesetting engine.

\end{document}
